%-------------------------------------------------------------------------------
%	SECTION TITLE
%-------------------------------------------------------------------------------
\iflanguage{german}{
	\cvsection{Über Mich}
}
{
	\cvsection{About Me}
}


%-------------------------------------------------------------------------------
%	CONTENT
%-------------------------------------------------------------------------------
\begin{cvparagraph}

\iflanguage{german}
{
	Ich bin ein technischer Agile Coach mit einem Software-Engineering-Hintergrund. Ich entwickle nach wie vor selbst Software. Mein Ansatz verbindet tiefes theoretisches Verständnis mit jahrelanger praktischer Erfahrung in Unternehmensumgebungen, in denen ich eng mit Teams und Führungskräften zusammenarbeite, um effektive, effiziente und stabile Systeme zu schaffen. Frameworks und Methoden ordne ich diesem Ziel unter.

	Ich leite die Entwicklung von \underline{\href{https://github.com/LetPeopleWork/Lighthouse}{Lighthouse}}, einem Open-Source-Tool für Flow-Metriken und probabilistische Prognosen (C\#/.NET-Backend, React-/TypeScript-Frontend), das von Dutzenden Teams in mehreren Unternehmen genutzt wird, inklusive zahlender Kunden. Releases erfolgen kontinuierlich im wöchentlichen Rhythmus, und ich wende dabei täglich TDD, CI/CD, Pair-/Mob-Programming sowie KI-unterstütztes Entwickeln an. Automatisierung ist eine weitere Leidenschaft. Ob ein Skript zur automatischen Datenerfassung oder \underline{\href{https://github.com/huserben/cv/actions}{die CI/CD-Pipeline hinter diesem Lebenslauf}}: solche Dinge zu bauen macht mir einfach Spass.

	Mein technischer Hintergrund ermöglicht es mir, Entwickler glaubwürdig zu coachen. Ob Teams eXtreme-Programming-Praktiken (TDD, Pair-/Mob-Programming, KI-unterstütztes Entwickeln) einführen oder ihre CI/CD- und DevOps-Fähigkeiten weiterentwickeln: ich habe selbst praktische Erfahrung darin, solche Veränderungen umzusetzen, nicht nur sie zu unterrichten.

	In den letzten Jahren habe ich mich verstärkt darauf konzentriert, datenbasierte Entscheidungsfindung zu fördern: durch den Einsatz von Flow-Metriken zur Messung und Verbesserung der Systemstabilität sowie probabilistischer Prognosen zur Vorhersage von Lieferterminen. Die Arbeit meiner Teams in diesen Bereichen wurde international anerkannt; wir haben Fallstudien und Workshops auf Konferenzen und Meetups weltweit vorgestellt.
}
{
	I'm a technical agile coach with a software engineering background, and I still ship code. My approach pairs deep theoretical grounding with years of hands-on experience in corporate environments, partnering with teams and leadership to build effective, efficient, and predictable systems regardless of framework or method.

	I lead development of \underline{\href{https://github.com/LetPeopleWork/Lighthouse}{Lighthouse}}, an open-source flow-metrics and probabilistic forecasting tool (C\#/.NET backend, React + TypeScript frontend) used by dozens of teams across multiple companies, including paying customers. Releases ship on a continuous, weekly cadence, and I practise TDD, CI/CD, pair/mob programming, and AI-augmented development day-to-day. Automation is another passion. From a data-collection script to \underline{\href{https://github.com/huserben/cv/actions}{the CI/CD pipeline behind this CV}}, I genuinely enjoy building these things.

	My engineering background lets me coach developers credibly. Whether teams are adopting eXtreme Programming practices (TDD, pair/mob programming, AI-augmented development) or evolving their CI/CD and broader DevOps capabilities, I have hands-on experience driving these changes, not just teaching them.

	In recent years I've focused on enabling data-informed decision-making: applying Flow Metrics to measure and improve system stability, and leveraging probabilistic forecasting to predict delivery timelines. My teams' work in these areas has gained international recognition, with case studies and workshops presented at conferences and meetups worldwide.
}
\end{cvparagraph}
